% ৪.১ Machine Language
\section{Machine Language}

\begin{learningbox}
এই টপিক শেষে শিক্ষার্থী machine language-এর সংজ্ঞা, বৈশিষ্ট্য, সুবিধা ও অসুবিধা ব্যাখ্যা করতে পারবে।
\end{learningbox}

\begin{definitionbox}{Machine language}
কম্পিউটারের processor সরাসরি যে binary instruction বা 0 ও 1-এর ভাষা বুঝতে পারে, তাকে machine language বলা হয়।
\end{definitionbox}

\begin{examplebox}{ধারণা}
\texttt{10110000 01100001} ধরনের binary pattern processor-এর কাছে instruction/data হিসেবে কাজ করতে পারে। মানুষ এটি সহজে বুঝতে পারে না, কিন্তু machine সরাসরি বুঝতে পারে।
\end{examplebox}

\begin{center}
\begin{tabular}{|p{0.30\textwidth}|p{0.50\textwidth}|}
\hline
\textbf{সুবিধা} & দ্রুত execute হয়; translator লাগে না। \\
\hline
\textbf{অসুবিধা} & লেখা, মনে রাখা, debug ও maintain করা কঠিন; machine dependent। \\
\hline
\end{tabular}
\end{center}

\begin{mistakebox}{সাধারণ ভুল}
Machine language দ্রুত হলেও programmer-friendly নয়। দ্রুত execution আর সহজ programming একই জিনিস নয়।
\end{mistakebox}
